\label{appendix:ensemble_tuning}

\rev{The ensemble size $n$ is a critical hyperparameter of the EnKF: too few members under-sample the prior and introduce spurious correlations, while an unnecessarily large ensemble increases computational cost without meaningful accuracy gains. To select an appropriate ensemble size for each cell line group, a systematic convergence study was performed by running the full dual estimation for every dataset at ensemble sizes $n \in \{10, 25, 50, 75\}$ and evaluating two complementary error metrics alongside wall-clock runtime.}

\rev{The \textbf{overall normalised RMSE} aggregates root-mean-square errors across all $n_x$ state variables and $n_r$ repeated runs, normalising each state by its median RMSE to prevent high-concentration species from dominating the score:}
\rev{
\begin{equation}
\overline{\mathrm{RMSE}}_{\mathrm{norm}} = \frac{1}{n_r}\sum_{r=1}^{n_r}\frac{1}{n_x}\sum_{j=1}^{n_x}\frac{\mathrm{RMSE}_{r,j}}{\widetilde{\mathrm{RMSE}}_j}
\label{eq:norm_rmse}
\end{equation}
where $\mathrm{RMSE}_{r,j}$ is the RMSE of the $j$th state in the $r$th run and $\widetilde{\mathrm{RMSE}}_j = \mathrm{median}_r(\mathrm{RMSE}_{r,j})$ is the median RMSE of state $j$ across all runs and ensemble sizes. A lower value indicates better overall accuracy.}

\rev{The \textbf{standard deviation of normalised RMSE} quantifies run-to-run variability:}
\rev{
\begin{equation}
\sigma_{\mathrm{norm}} = \mathrm{std}_{r}\!\left(\frac{1}{n_x}\sum_{j=1}^{n_x}\frac{\mathrm{RMSE}_{r,j}}{\widetilde{\mathrm{RMSE}}_j}\right)
\label{eq:std_rmse}
\end{equation}
A lower value indicates more reproducible estimation across independent ensemble initialisations.}

\rev{Figure~\ref{fig:ensemble_tuning} summarises the results. For the three \textbf{Cell Line~A} datasets (panels a--c), the overall normalised RMSE reaches a minimum at $n = 50$; increasing to $n = 75$ yields no statistically meaningful improvement while the runtime approximately doubles. An ensemble size of 50 was therefore selected for all Cell Line~A conditions.}

\rev{For \textbf{Cell Line~B}, which differs from the base model calibration cell line in host background, product, and feeding strategy, the picture is more nuanced. For Feed~C (panel~d) and Feed~U (panel~e), the normalised RMSE continues to decrease from $n = 50$ to $n = 75$, indicating that the larger ensemble provides additional benefit when the model--data mismatch is substantial. For Feed~U~+~40\% (panel~f), $n = 50$ and $n = 75$ show comparable accuracy; nevertheless, $n = 75$ was adopted uniformly across all three Cell Line~B datasets to ensure a consistent comparison across feeding strategies and to provide a more conservative uncertainty representation when transferring to a new cell line.}

\begin{figure}[H]
\centering
\includegraphics[width=\textwidth]{Figs/ensemble_tuning.png}
\caption{\rev{Ensemble size selection for Cell Line~A (a--c) and Cell Line~B (d--f). Each panel shows the overall normalised RMSE (red, left axis), the standard deviation of normalised RMSE (blue, left axis), and the wall-clock runtime for the entire culture simulation (green, right axis) as a function of ensemble size $n \in \{10, 25, 50, 75\}$.}}
\label{fig:ensemble_tuning}
\end{figure}

\rev{\textbf{Computational cost.} All simulations were performed in Python~3.14 on a desktop workstation equipped with a 13th-generation Intel Core i5-13600K processor (14 cores, 20 threads, 3.5~GHz base clock) and 32~GB DDR4 RAM, running Windows~11 (64-bit). The mechanistic model was integrated using SciPy's explicit Runge--Kutta solver (\texttt{solve\_ivp}, RK45) with an internal time step of 0.01~h ($\approx$36~s). With the selected ensemble sizes ($n = 50$ for Cell Line~A, $n = 75$ for Cell Line~B), a complete simulation of the full fed-batch culture (12--14 days, 7--8 measurement updates) required approximately 2--3 minutes of wall-clock time. This runtime is well within the timescales of practical interest for fed-batch monitoring, where off-line measurements are typically available every 24~h, confirming the feasibility of the framework for near-real-time state estimation and decision support in manufacturing environments.}
